We compute all six metrics for every district produced by the four \textsc{bisect}
structure algorithms on NC ($k=14$), WI ($k=8$), and TX ($k=38$) under 2020 Census data,
using single seed $= 42^\dagger$.
This yields $4 \times (14 + 8 + 38) = 240$ district observations per metric,
for a total of $1{,}440$ (metric, district) pairs.

\subsection{Within-Family Correlations}

Table~\ref{tab:corr} reports the $6\times 6$ Pearson correlation matrix
aggregated across all 240 districts and all four algorithms.

\begin{table}[ht]
\centering
\caption{Pairwise Pearson correlation among six compactness metrics
across 240 districts (NC+WI+TX, all four structure algorithms, seed $= 42^\dagger$).}
\label{tab:corr}
\begin{tabular}{lrrrrrr}
\toprule
       & PP    & Reock & CH    & S     & LW    & PWC \\
\midrule
PP     & 1.00  & 0.71  & 0.58  & $-$0.99 & $-$0.62 & 0.34 \\
Reock  & 0.71  & 1.00  & 0.79  & $-$0.70 & $-$0.71 & 0.31 \\
CH     & 0.58  & 0.79  & 1.00  & $-$0.58 & $-$0.55 & 0.27 \\
S      & $-$0.99 & $-$0.70 & $-$0.58 & 1.00 & 0.62  & $-$0.34 \\
LW     & $-$0.62 & $-$0.71 & $-$0.55 & 0.62 & 1.00  & $-$0.29 \\
PWC    & 0.34  & 0.31  & 0.27  & $-$0.34 & $-$0.29 & 1.00 \\
\bottomrule
\end{tabular}
\begin{tablenotes}
\footnotesize
\item Note: Negative correlations with S and LW reflect their inverted scales
  (higher S = less compact; higher LW = more elongated).
\end{tablenotes}
\end{table}

Key findings from Table~\ref{tab:corr}:

\paragraph{Within-family: area-based.}
Reock and CH correlate at $r = 0.79$, confirming they measure related geometric
properties. Both capture departure from convexity or circularity at the
``envelope'' level.

\paragraph{Within-family: perimeter-based.}
PP and S correlate at $r = -0.99$ (negative because S increases when PP decreases).
This near-perfect correlation confirms the $S = 1/\sqrt{\text{PP}}$ identity:
after accounting for the monotone transformation, these metrics carry
virtually identical information.

\paragraph{Across families.}
PP and Reock correlate at $r = 0.71$---the strongest cross-family correlation
(documented in K.1 and K.2). PP and CH correlate at $r = 0.58$.
These moderate values ($< 0.80$) confirm that perimeter-based and area-based
metrics are related but non-redundant.

\paragraph{PWC orthogonality.}
PWC correlates at $r \leq 0.34$ with all geometric metrics.
This confirms that representational compactness (PWC) and geometric compactness
(all other metrics) measure independent properties, consistent with
\citet{fryer2011measuring}.

\subsection{State-Level Variation}

Correlation structure varies modestly by state.
In TX ($k = 38$), the PP--Reock correlation drops to $0.64$ due to the higher
prevalence of mixed urban-rural districts where the MBC is determined by
distant rural tracts while the perimeter is driven by urban boundary detail.
In WI ($k = 8$), the Reock--CH correlation rises to $0.83$ because WI's
convex state boundary means that border districts have convex hulls close
to their MBC geometry.

\subsection{Algorithm-Level Variation}

The moving-knife algorithm (MKA) consistently produces higher Reock values
and lower LW values than other algorithms, consistent with its design goal
of minimising the worst-case Reock score (K.2, Claim 2).
Prime-factor produces lower PWC values than standard-bisect across all three states
(K.6, Claim 2), driven by its hierarchical nesting aligning population centres
with district centroids.
